\documentclass[10pt]{article}
\usepackage[utf8]{inputenc}
\usepackage[T1]{fontenc}
\usepackage{amsmath}
\usepackage{amsfonts}
\usepackage{amssymb}
\usepackage[version=4]{mhchem}
\usepackage{stmaryrd}
\usepackage{bbold}

%New command to display footnote whose markers will always be hidden
\let\svthefootnote\thefootnote
\newcommand\blfootnotetext[1]{%
  \let\thefootnote\relax\footnote{#1}%
  \addtocounter{footnote}{-1}%
  \let\thefootnote\svthefootnote%
}

%Overriding the \footnotetext command to hide the marker if its value is `0`
\let\svfootnotetext\footnotetext
\renewcommand\footnotetext[2][?]{%
  \if\relax#1\relax%
    \ifnum\value{footnote}=0\blfootnotetext{#2}\else\svfootnotetext{#2}\fi%
  \else%
    \if?#1\ifnum\value{footnote}=0\blfootnotetext{#2}\else\svfootnotetext{#2}\fi%
    \else\svfootnotetext[#1]{#2}\fi%
  \fi
}

\begin{document}
This is a guide to how to think about the midterm. It is not a literal practice exam but mastering these concepts will be a good preparation for the exam. The midterm exam will be closed book with no calculators. The exam will be $\mathbf{7 5}$ minutes. The midterm will cover up to and including the material in Section 5.2 of the textbook.

\section*{1. Key Topics}
Topics of focus include (but are not limited to):\\
(1) Arithmetic Functions. Some key examples of arithmetic functions include: Möbius Function $\mu(n)$, Euler $\phi$-function $\phi(n)$, number divisors function $v(n)$, sum of divisors function $\sigma(n)$.\\
(2) Multiplicative and completely multiplicative functions.\\
(3) Convolution of arithmetic functions and Möbius inversion\\
(4) Quadratic residues and Legendre symbols\\
(5) Evaluation of special Legendre symbols e.g. $\left(\frac{-1}{p}\right)$ and $\left(\frac{2}{p}\right)$ for $p$ prime\\
(6) Quadratic Reciprocity

$$
\left(\frac{p}{q}\right)\left(\frac{q}{p}\right)=(-1)^{\frac{(p-1)}{2} \cdot \frac{(q-1)}{2}} \quad \text { for } \quad p \neq q \quad \text { odd primes. }
$$

(7) Definition of primitive roots. Definition of the order of an integer modulo $m$. (The Primitive root theorem will not be covered by this midterm).\\
Note: Although the topics above are the main focus, you will still need to know the material covered by the first midterm to fully understand the above material.

\section*{2. Practice Questions}
Question 1. Computational questions:

\begin{itemize}
  \item Is $x^{2} \equiv 137(\bmod 401)$ solvable? (note: 137 and 401 are both prime).
  \item Find all incongruent quadratic residues modulo 13. Do the same for the modulus 17, and the modulus 23.
  \item Find all incongruent solutions to the congruence $x^{2} \equiv 23(\bmod 77)$.
  \item Find all incongruent solutions to the congruence $x^{2} \equiv 11(\bmod 39)$. [HINT: For the above two questions use the Chinese Remainder Theorem].
  \item Prove that von Mangoldt function
\end{itemize}

$$
\Lambda(n):= \begin{cases}\log p & n=p^{a} \text { such that } p \text { prime and } a \in \mathbb{N} \\ 0 & \text { otherwise }\end{cases}
$$

is not multiplicative.

\footnotetext{Date: Fall 2023.
}\begin{itemize}
  \item Use Euler's criterion to determine $\left(\frac{-6}{11}\right)$.
  \item Use Gauss' Lemma to determine $\left(\frac{12}{23}\right)$.
  \item Use Quadratic Reciprocity to determine ( $\frac{91}{127}$ ).
  \item Find the order of 8 modulo 15. Find the order of 9 modulo 17.
  \item Assuming there is a primitive root modulo 18 , how many incongruent primitive roots are there? Do the same question for modulus 19 and modulo 57.
  \item What are all the possible orders of an integer modulo 31.
  \item Compute the order of 2 modulo $1085=5 \cdot 7 \cdot 31$ (this is not meant to be done by brute force).
  \item Let $p$ and $q=4 p+1$ be primes. Compute $\left(\frac{p}{q}\right)$.
  \item If $a$ is a primitive root modulo 31 , what is the order of $a^{3}, a^{12}, a^{16}, a^{17}$, $a^{20}$.
\end{itemize}

Question 2. Some Theorems to know by name and how they can be applied to solve problems.

\begin{itemize}
  \item Euler's Criteria for the Legendre symbol (and proof)
  \item Gauss' Lemma for the Legendre Symbol (and proof)
  \item Möbius Inversion
  \item Quadratic Reciprocity
\end{itemize}

Question 3. These questions require the use of one of the named theorems above. Correct solutions should cite the theorem used and explain the computation/steps.

\begin{itemize}
  \item Show that $\sum_{a=1}^{p-1}\left(\frac{a}{p}\right)=0$.
  \item Suppose that $p$ and $2 p+1$ are both odd primes. Find all incongruent solutions to
\end{itemize}

$$
x^{2} \equiv 1 \quad(\bmod p(2 p+1))
$$

\begin{itemize}
  \item Let $p$ be an odd prime and let a be a quadratic residue mod $p$. Show that if $-a$ is also a quadratic residue modulo $p$, then $p \equiv 1(\bmod 4)$.
  \item Derive a closed form formula for $\left(\frac{5}{p}\right)$ for primes $p \geq 7$.\\[0pt]
[HINT: The answer should be a piecewise function whose values depend on some congruence classes for the prime $p$ modulo 5].
  \item If a is a quadratic residue modulo two distinct odd primes $p$ and $q$, then the equation $x^{2} \equiv a(\bmod p q)$ is solvable.
  \item Use Gauss' Lemma to argue that $\left(\frac{2}{p}\right)=(-1)^{N}$, where $N=\frac{p-1}{2}-\left\lfloor\frac{p}{4}\right\rfloor$.
\end{itemize}

Question 4. - Define the Möbius function $\mu(n)$. Show that $\mu$ it is multiplicative, but not completely multiplicative.

\begin{itemize}
  \item Let $f$ be a multiplicative arithmetic function. Show that $F(n)=\sum_{d \mid n} f(d)$ is multiplicative.
  \item Show that $\sum_{d \mid n} \phi(d)=n$ for all integers $n \geq 1$.
  \item Evaluate $\sum_{d \mid n, d>0} \mu(d)$ for all integers $n \geq 1$.
  \item Evaluate $\sum_{d \mid n, d>0} \mu(d) v(n / d)$, where $v$ is the number of divisors function and $n=4^{3} \cdot 7^{2}$.
  \item Show that $\sum_{d \mid n, d>0} \frac{1}{d}=\frac{\sigma(n)}{n}$ for all integers $n \geq 1$.
  \item For integers $k \geq 1$, and multiplicative function $f$, show that the function $g(n):=f\left(n^{k}\right)$ is multiplicative.
  \item Evaluate the number divisor function $v(450), v(105), v\left(p^{k}\right)$, for prime $p$ and non-negative integer $k$.
\end{itemize}


\end{document}